\subsection{Excluding $C_{66}$}
\begin{lemma}\label{lem:decic-local-two-cycle}
Let $K$ be a number field, let $\sigma\in\operatorname{Aut}(K)$ fix a
prime $\mathfrak p$ above $2$, and suppose that $E/K$ has ordinary good
reduction at $\mathfrak p$ and $\operatorname{End}_{\overline K}(E)=\mathbb Z$.
There is no $K$-rational isogeny of degree $2$ from $E$ to $E^\sigma$.
\end{lemma}
\begin{proof}
Put $h=\operatorname{ord}(\sigma)$ and $e=v_{\mathfrak p}(2)$.
The composite of the $h$ conjugate isogenies is an endomorphism of $E$
of degree $2^h$. Thus $h$ is even and the composite is
$[\pm2^{h/2}]$. With respect to minimal N\'eron differentials, a degree-$2$
isogeny between ordinary elliptic curves has differential multiplier
of valuation either $0$ or $e$, according as its kernel is \'etale or
connected. Since $\sigma$ fixes $\mathfrak p$, all $h$ conjugate
multipliers have the same valuation. Their product therefore has
valuation $0$ or $he$, whereas $[\pm2^{h/2}]$ has valuation $he/2$.
This is a contradiction.
\end{proof}

\begin{proposition}\label{prop:decic66-exclusion}
An elliptic curve over a number field of degree $10$ has no point of order $66$.
\end{proposition}
\begin{proof}
Put $X=X_1(33)$ and $J=J_1(33)$.
Suppose that $E/K$ has a point of order $66$, where $[K:\mathbb Q]=10$.
The point of order $33$ gives a rational effective divisor $D$ of degree
$10$ on $X$, and the point of order $2$ gives a $K$-rational isogeny
$\phi:E\to E'$. Let $D'$ be the divisor obtained from $E'$ and the image
of the marked $33$-torsion point.

By \cite[Theorem~3.1(2), Theorem~4.13 and Table~2]{DEHMZ},
\[
 J(\mathbb Q)\simeq C_5\times C_{42373650}
\]
is generated by rational cuspidal divisors.
The retained proven class-group computation, repeated in the new
packet, gives
\[
 J(\mathbb F_2)\simeq C_{15}\times C_{42373650}.
\]
The global cuspidal generators have image
$H\simeq C_5\times C_{42373650}$, so reduction is injective on the
whole of $J(\mathbb Q)$ and has image $H$.
We write $\chi:J(\mathbb F_2)\to C_3$ for the quotient by $H$.
The global closed cusp degrees are $1^{10},2^5,10^2$, and the finite
cusp degrees are $1^{10},2^5,5^2,10$. All rational and quadratic
cusps lift intact. The two degree-$5$ finite cusps are the two parts
of one global degree-$10$ cusp. Their $\chi$-values are nonzero and
opposite. The signatures $(v_c(b^4/\Delta),v_c(b))$ identify the
complete global cusp blocks in both good characteristics used below.

We first describe the degree-$10$ reduction classes.
The exact noncuspidal place counts over $\mathbb F_2$ are
\[
 (0,0,0,0,2,20,30,25,50,123).
\]
Among the places of degree $d$, the numbers whose degree-zero class
lies in $H$ are
\[
\begin{array}{c|rrrrrr}
 d&5&6&7&8&9&10\\\hline
 \#&0&0&10&5&20&21.
\end{array}
\]
Every effective divisor of degree $10$ has at most two noncuspidal
components, and two are possible only in degrees $5+5$.
The complete enumeration retains all cusp multiplicities.
There are $321230$ effective divisors whose classes lie in the rational
image, in $321210$ classes. Exactly ten classes have $h^0=2$;
each contains three effective divisors, all cuspidal. All other
classes have $h^0=1$. In particular, every retained divisor containing
a noncuspidal component has $h^0=1$.
These dimensions follow also by counting the effective divisors in
each class over $\mathbb F_2$, where their number is $2^{h^0}-1$.
A separate generating-function convolution verifies the complete count.
The exact Frobenius polynomial verifies completeness of the place
list, and the independent finite-field replay binds every one of the
$250$ places to an actual marked elliptic curve.

If a degree-$5$ noncuspidal component occurs with multiplicity one,
$E$ has supersingular good reduction at an unramified prime above $2$.
This is incompatible with a rational subgroup of order $2$: the
isogeny and its dual both have inseparable special fibres, so both
N\'eron differential multipliers are divisible by $2$, whereas their
product is the multiplier of $[2]$ and has valuation $1$.
The double of either degree-$5$ place has nonzero $\chi$-value and is
already excluded. This also treats the divisor consisting of the two
distinct degree-$5$ places. Degree-$6$ components are excluded by
$\chi$, since their cusp complements have degree $4$ and lift globally.

For degrees $7$ and $8$, the retained places are reductions of the ten
known septic and five known octic divisors respectively. The accepted
witness coordinates and their Picard classes are unchanged in the new
frame. After adding the global cusp complement, injectivity of
specialization identifies the class of $D$ with that known effective
divisor. Since its special fibre has $h^0=1$, semicontinuity identifies
the divisors themselves, contradicting the noncuspidal support of $D$.
Of the twenty degree-$9$ classes, ten are reductions of the known
nonic divisors and give the same contradiction after adding a rational
cusp. The other ten are the signed global cuspidal classes $C_i$ of
degree $9$ in the retained $X_1(33)$ packet. For each of the ten rational
cusps $c$, the new computation gives
\[
 h^0(X_{\mathbb F_5},\overline{C_i+c})=0.
\]
Thus these $100$ classes have no effective global representatives.
For an independent certificate, we use a common effective pole divisor
$A$ of degree $58$. The packet supplies a full basis of the
$38$-dimensional space $L(A)$, and all $100$ restrictions to
$A-C_i-c$, of degree $48$, have rank $38$. A fresh native reconstruction
checks all function poles and $72$ local prefix kernels; independent
coefficient algebra and Gaussian elimination verify all ranks.

Suppose next that $\overline D$ is a noncuspidal closed point of degree
$10$. There is a unique prime of $K$ above $2$, unramified of residue
degree $10$, and $E$ has good reduction there. The supersingular case
is excluded as above. In the ordinary case, the reduction of $D'$ is
$\overline D$ or $\langle2\rangle\overline D$, according as $\phi$ is
Frobenius or Verschiebung. Injectivity of $J(\mathbb Q)$ specialization
gives the corresponding equality of global divisor classes.
The special class has $h^0=1$, so
\[
 D'=D\quad\text{or}\quad D'=\langle2\rangle D.
\]
Fine moduli then gives a $K$-rational $2$-isogeny from $E$ to $E^\sigma$
for some $\sigma\in\operatorname{Aut}(K)$. The prime above $2$ is fixed
by $\sigma$. By \cite[Section~4.10]{CCRS}, $E$ is not geometrically CM,
since no decic CM torsion group contains $C_{33}$.
Lemma~\ref{lem:decic-local-two-cycle} gives a contradiction.

It remains that $\overline D$ is entirely cuspidal. Every cuspidal
reduction class in the rational image has a global effective cuspidal
representative. If its $h^0$ is $1$, injectivity and semicontinuity
give the same contradiction as before. We are therefore left with
ten global line bundles of degree $10$ and dimension $h^0=2$.
Indeed, each has three distinct global effective cuspidal
representatives with the same reduction class. Injectivity makes
them linearly equivalent globally, and the finite dimension $2$
is an upper bound. By \cite[Table~1]{DvH},
$\operatorname{gon}_{\mathbb Q}(X)=10$, so these pencils are base point free.
They form two diamond orbits of five pencils.
The same argument applies to $D'$, since an isogeny preserves
potentially multiplicative reduction. Thus $D$ and $D'$ both lie in
this set of ten pencils.

We now compare these pencils at $5$. The dimensions remain $2$ there.
The chosen reduced functions all have degree $10$, so the special
pencils are base point free as well.
Consequently their global section lattices specialize onto the full
spaces of sections at $5$, so each pencil has a morphism to
$\mathbb P^1_{\mathbb Z_5}$ with the indicated special fibre.
Choose representatives $f_A,f_B$ of the two orbits.
They are constructed by Riemann--Roch from complete global cusp blocks,
and all ten diamond translates are identified by their exact
cuspidal pole classes. The following table describes their fibres
over $\mathbb P^1(\mathbb F_5)$:
\[
\begin{array}{c|c|c|c}
 &\text{entirely cuspidal fibres}&
 \text{other fibre degrees on }X&
 \text{maximum common degree after a marked }2\text{-isogeny}\\\hline
 f_A&0,1,\infty&(4+6),10,10&4\\
 f_B&0,4,\infty&(2+8),(3+3+4),(2+8)&4.
\end{array}
\]
The last column ranges over every target pencil and every rational
target value, and counts multiplicities. More precisely, all nonzero
bounds occur for source value $2$, and are $4$; the other noncuspidal
source fibres give bound $0$.

For clarity, this is a computation on the marked $2$-isogeny curve
$X_1(66)$. Over a noncuspidal place of degree $d$ on $X$, factor the
$2$-division cubic. A factor of degree $e$ gives a place of degree
$de$ on $X_1(66)$; the inverse marking is
$P_{66}=17P_{33}+R_2$. The fibre multiplicity is retained, so its
contribution is $mde$. There are $23$ kernel orbits above the eleven
noncuspidal places in the six relevant source fibres. We compute all
marked quotient curves and every target-pencil value. Independent
Sage V\'elu maps verify the actual markings, all source field degrees,
the saved prime generators, and every function value. Thus the last
column bounds the degree of the actual common fibre, including all
residue norms.

Apply a simultaneous diamond so that the source pencil is $f_A$ or
$f_B$. Both its value and the value of the target pencil are rational,
so their reductions belong to $\mathbb P^1(\mathbb F_5)$.
If the source fibre is noncuspidal, the rational effective divisor of
degree $10$ on $X_1(66)$ is contained in a common fibre of degree at
most $4$, a contradiction. Otherwise it reduces entirely to cusps.
This is also impossible. Indeed, $J_1(66)(\mathbb Q)$ has rank zero
by \cite[Theorem~3.1(2)]{DEHMZ}, and
\[
 \operatorname{gon}_{\mathbb Q}(X_1(66))
 \geq \frac{7\cdot1440}{800}>10.
\]
At the odd good prime $5$, lift the cuspidal reduction to a
$\mathbb Q_5$-rational cuspidal divisor $C$ of degree $10$.
The class of its difference from the given global divisor is torsion
by rank zero and Manin--Drinfeld, and reduces to zero. The formal
kernel over $\mathbb Q_5$ is torsion-free, so the two divisors are
linearly equivalent. Hence the global degree-$10$ line bundle has at
least two sections after extension to $\mathbb Q_5$, and therefore
over $\mathbb Q$, contradicting the gonality bound. This completes
the proof.
\end{proof}
